\chapter[1929 --- Eijkman et al.]{1929: Christiaan Eijkman and Sir Frederick Gowland Hopkins}
\label{ch:1929_Eijkman_Hopkins}

\section*{Main Contribution}

\textit{Established that essential dietary factors are required for health and growth, helping lay the foundation for vitamin science.}

\section{Official Citation}

for their discoveries of vitamins

\section{Background and Historical Context}

Severe deficiency diseases had long been associated with particular diets, but the responsible factors were not understood. Research on beriberi and animal growth provided evidence that food contains essential substances beyond proteins, fats, carbohydrates, and minerals.

\section{The Discovery and Contribution}

Christiaan Eijkman's studies of beriberi in Java showed that polished rice and unpolished rice had different effects on disease in experimental animals. His observations supported the existence of an antineuritic factor associated with the rice husk and helped establish the nutritional origin of beriberi.

Frederick Gowland Hopkins demonstrated that animals require additional growth-promoting substances in their diet. His work provided a broader biochemical framework for the idea that small, essential dietary factors are necessary for normal growth and health.

The Nobel Prize in Physiology or Medicine 1929 was divided equally between Eijkman, ``for his discovery of the antineuritic vitamin,'' and Hopkins, ``for his discovery of the growth-stimulating vitamins.''

\section{Impact on Modern Medicine}

Their discoveries helped establish nutritional deficiency as a preventable cause of disease and supported the later identification, purification, and clinical use of individual vitamins. Vitamin supplementation, food fortification, and nutritional epidemiology all grew from this foundation.

\section{Legacy}

Eijkman and Hopkins helped transform nutrition from an empirical field into a biochemical science. Their work continues to inform the prevention and treatment of deficiency diseases worldwide.

\section{Modern Developments}

Modern nutrition research examines vitamin metabolism, genetic differences in nutrient requirements, food fortification, and the relationship between micronutrient status and chronic disease.
